%! AP Statistics — Unit 1, Lesson 1.3 — Lesson Plan
\documentclass[10pt]{article}
\usepackage[margin=0.6in]{geometry}
\usepackage{apstats-boxes}
\usepackage{enumitem}
\usepackage{multicol}
\usepackage{amsmath,amssymb}
\usepackage{array}
\usepackage{tabularx}
\usepackage{booktabs}
\usepackage{colortbl}

\newcolumntype{C}[1]{>{\centering\arraybackslash}p{#1}}
\newcolumntype{Y}{>{\centering\arraybackslash}X}
\setlist[itemize]{leftmargin=1.2em,itemsep=0.35em,topsep=0.35em}
\setlist[enumerate]{leftmargin=1.4em,itemsep=0.35em,topsep=0.35em}
\setlength{\parindent}{0pt}
\setlength{\parskip}{3pt}

\begin{document}

\begin{center}
    {\Large\bfseries AP Statistics: 2026--2027\ (55 minutes) \\[4pt]
    {\small\bfseries Unit 1: Exploring One-Variable Data \qquad
     Lesson 1.3: Representing a Categorical Variable with Tables}}
\end{center}

\begin{tcolorbox}[colback=sky,colframe=navy,boxrule=0.9pt,arc=2mm,
    left=2mm,right=2mm,top=1.5mm,bottom=1.5mm]
    \textbf{Primary Objective:} Students will construct and interpret frequency and
    relative frequency tables for a categorical variable. Students will recognize that
    relative frequency (not raw count) is necessary when comparing groups of different
    sizes, and will use tables to describe the distribution of a categorical variable
    in context (Big Idea: UNC).
\end{tcolorbox}

\vspace{4pt}

\begin{skillbox}[Priority Ideas \& Skills]{goldbox}
\begin{minipage}[t]{0.48\linewidth}
    \textbf{AP Skill 2 --- Data Analysis}
    \begin{itemize}
        \item Construct a frequency table and a relative frequency table from raw
              categorical data.
        \item Interpret relative frequencies in context
              (e.g., ``32\% of students prefer\ldots'').
        \item Explain why relative frequency is used to compare groups of different sizes.
    \end{itemize}
\end{minipage}
\hfill
\begin{minipage}[t]{0.48\linewidth}
    \textbf{Key Understandings}
    \begin{itemize}
        \item Raw counts are misleading when total group sizes differ.
        \item Relative frequencies must sum to 1 (or 100\%) --- this is a built-in
              error check.
        \item A frequency table is an alternative to a bar chart for displaying the
              distribution of a categorical variable; use one when precision matters.
    \end{itemize}
\end{minipage}
\end{skillbox}

\begin{skillbox}[{Vocabulary, Concepts, \& Theorems}]{greenbox}
\begin{minipage}[t]{0.48\linewidth}
    \begin{itemize}
        \item \textbf{Frequency} --- the count of individuals in a category
        \item \textbf{Relative frequency} --- frequency $\div$ total; expressed as a
              proportion or percentage
        \item \textbf{Frequency table} --- lists each category with its count
        \item \textbf{Relative frequency table} --- lists each category with its
              relative frequency
    \end{itemize}
\end{minipage}
\hfill
\begin{minipage}[t]{0.48\linewidth}
    \begin{itemize}
        \item \textbf{Distribution (categorical)} --- the set of possible values and
              their associated relative frequencies
        \item \textbf{Proportion} --- a relative frequency expressed as a decimal
        \item \textbf{Percent} --- a relative frequency expressed per 100
        \item \textbf{Mode} --- the category with the highest frequency (note: rarely
              the focus on AP exams)
    \end{itemize}
\end{minipage}
\end{skillbox}

\begin{skillbox}[Activate Prior Knowledge \& Spiral Review]{sky}
\begin{minipage}[t]{0.48\linewidth}
    \textbf{Warm-Up Worksheet (5 min)}\\
    Students receive a list of 20 students' favorite school subjects (categorical
    data). Tasks: tally each category; ask ``Can we tell which subject is most
    popular just by scanning the list?'' Motivates today's tables as an
    organizational tool.\\[6pt]
    \textbf{Topics Connected}
    \begin{itemize}
        \item[\rule{0.7cm}{0.4pt}] Categorical variable (Lesson 1.2)
        \item[\rule{0.7cm}{0.4pt}] Distribution of a variable (Lesson 1.2 preview)
        \item[\rule{0.7cm}{0.4pt}] Preview: displaying data with bar charts (Lesson 1.4)
    \end{itemize}
\end{minipage}
\hfill
\begin{minipage}[t]{0.48\linewidth}
    \textbf{Connections to Current Learning}
    \begin{itemize}
        \item Lesson 1.4 displays the same categorical distributions \textit{visually}
              (bar charts, pie charts). Tables and graphs answer the same question
              differently.
        \item Relative frequency now is the same operation as computing a probability
              in Unit 4.
        \item Two-way frequency tables (Unit 2) are a direct extension of today's
              one-way tables.
    \end{itemize}
\end{minipage}
\end{skillbox}

\begin{skillbox}[Hook \& Lesson]{sky}
\begin{minipage}[t]{0.48\linewidth}
    \textbf{Hook (4--5 min)}\\
    Display counts for two schools' survey results on commute mode:
    \begin{center}
    \small\renewcommand{\arraystretch}{1.3}
    \begin{tabular}{lcc}
    \hline
    Mode & School A & School B \\ \hline
    Bus  & 120 & 30 \\
    Car  & 80  & 50 \\
    Walk & 40  & 120 \\ \hline
    \textbf{Total} & \textbf{240} & \textbf{200} \\
    \hline
    \end{tabular}
    \end{center}
    Ask: ``Which school has more walkers?'' Let students debate. Reveal: School B
    has \textit{fewer} walkers by count but \textit{more} by proportion. Transition
    to why we need relative frequency.
\end{minipage}
\hfill
\begin{minipage}[t]{0.48\linewidth}
    \textbf{Lesson Flow (15 min)}
    \begin{enumerate}
        \item From raw class data (birth month), build a tally, then a frequency
              table together.
        \item Add the relative frequency column: divide each count by $n$; verify
              sum $= 1$.
        \item Revisit the hook: compute relative frequencies for both schools and
              resolve the debate.
        \item Model AP-style interpretation:
              \textit{``About 60\% of School B students walk to school, compared
              to only 17\% of School A students.''}
    \end{enumerate}
\end{minipage}
\end{skillbox}

\begin{skillbox}[Explicit Instruction \& Active Monitoring]{sky}
\begin{minipage}[t]{0.48\linewidth}
    \textbf{Direct Instruction (10 min)}
    \begin{itemize}
        \item Model the four-step table build: \textbf{(1)} list categories,
              \textbf{(2)} tally, \textbf{(3)} count, \textbf{(4)} divide by $n$.
        \item Common errors: forgetting to divide (reporting counts as ``relative
              frequency''), rounding that makes sum $\neq 1$.
        \item AP language: always state the context and the total
              (\textit{``of the 200 students surveyed\ldots''}).
    \end{itemize}
\end{minipage}
\hfill
\begin{minipage}[t]{0.48\linewidth}
    \textbf{Active Monitoring}
    \begin{itemize}
        \item Circulate as students complete guided notes practice.
        \item Check: relative frequencies expressed as decimals \textit{and} as
              percentages (AP uses both; students should convert fluently).
        \item Prompt cold-calls: ``Does your sum equal 1? If not, where did the
              rounding error occur?''
        \item Use the notes guided-practice table as the monitoring anchor.
    \end{itemize}
\end{minipage}
\end{skillbox}

\begin{skillbox}[Group Work \& Differentiation]{redbox}
\begin{minipage}[t]{0.48\linewidth}
    \textbf{Group Activity (8--10 min)}\\
    Groups receive a raw dataset of 30 students' favorite sport (6 categories).
    Each group:
    \begin{enumerate}
        \item Tallies the data and builds a complete frequency table.
        \item Adds a relative frequency column; verifies the sum.
        \item Writes two sentences interpreting the distribution.
        \item A second dataset (different $n$) prompts comparison using relative
              frequency.
    \end{enumerate}
\end{minipage}
\hfill
\begin{minipage}[t]{0.48\linewidth}
    \textbf{Differentiation}
    \begin{itemize}
        \item \textit{Support}: Pre-labeled table template with categories already
              listed; student only fills counts and relative frequencies.
        \item \textit{Extension}: Two datasets of different sizes; compare
              distributions and write a short conclusion paragraph.
        \item \textit{AP Connection}: Discuss how an AP free-response item asks to
              ``describe the distribution'' --- complete sentences, context, and
              relative frequency are required.
    \end{itemize}
\end{minipage}
\end{skillbox}

\begin{skillbox}[Individual Work \& Assessment]{redbox}
\begin{minipage}[t]{0.48\linewidth}
    \textbf{Exit Ticket (5 min)}\\
    A survey of 50 students on preferred study location (library 18, home 22,
    caf\'{e} 6, other 4). Students complete the relative frequency column and write
    one sentence describing the most common preference in context.
\end{minipage}
\hfill
\begin{minipage}[t]{0.48\linewidth}
    \textbf{AP-Style Check}\\
    Free-response prompt: \textit{``A table shows the distribution of eye color for
    200 people. What proportion have brown or hazel eyes? Describe what the table
    tells you about the distribution of eye color in this sample.''}\\[4pt]
    Assess: correct proportion calculation, complete contextual sentence,
    identification of the most and least common category.
\end{minipage}
\end{skillbox}

\begin{skillbox}[Reinforcement \& Extension]{goldbox}
\begin{minipage}[t]{0.48\linewidth}
    \textbf{Homework / Practice}
    \begin{itemize}
        \item Given a raw list of 25 categorical values, construct a frequency and
              relative frequency table; describe the distribution in 2--3 sentences.
        \item Three AP-style problems: one table-build, one interpretation, one
              compare-two-groups using relative frequency.
    \end{itemize}
\end{minipage}
\hfill
\begin{minipage}[t]{0.48\linewidth}
    \textbf{Extension / Preview}
    \begin{itemize}
        \item \textit{Extension}: Find a published frequency table (U.S.\ Census,
              Pew Research) and write a description of the distribution.
        \item \textit{Preview}: Lesson 1.4 --- the same distributions displayed as
              bar charts and pie charts. Think: what does a table show that a bar
              chart doesn't (and vice versa)?
    \end{itemize}
\end{minipage}
\end{skillbox}

\end{document}
